% Research proposal, formatted to match CS498-researchProposal.pdf
%
% Self-contained: compiles with pdflatex and needs no external style file. If
% you have the NeurIPS style you used before, swap the preamble for
% \usepackage[final]{neurips_2024} and delete the manual title block below.

\documentclass[11pt]{article}

\usepackage[margin=1in]{geometry}
\usepackage[T1]{fontenc}
\usepackage[utf8]{inputenc}
\usepackage{microtype}
\usepackage{times}
\usepackage{booktabs}
\usepackage{array}
\usepackage{url}
\usepackage[hidelinks]{hyperref}
\usepackage{titlesec}

\titleformat{\section}{\large\bfseries}{\thesection}{0.8em}{}
\titlespacing*{\section}{0pt}{1.4em}{0.5em}

\setlength{\parindent}{0pt}
\setlength{\parskip}{0.55em}

\renewcommand{\abstractname}{Abstract}

\title{\textbf{Measuring Handoff Failure Between Independently\\ Trained Reinforcement Learning Skills}}
\author{%
  Isaiah Broderson\\
  University of Hawai\kern0.05em'\kern-0.05em i at Hilo\\
  \texttt{isaiahkb@hawaii.edu}
}
\date{}

\begin{document}
\maketitle

\begin{abstract}
\noindent
Long-horizon control problems are rarely solved by a single reinforcement learning
policy. The standard approach is to train narrow skills separately and execute them
in sequence, which replaces one hard credit-assignment problem with a set of easy
ones. This decomposition introduces a failure mode that is seldom measured directly:
each skill is trained on its own distribution of starting states, but when chained it
receives whatever state its predecessor produced. In this project I will train two
spacecraft control skills independently inside an existing physics simulator, chain
them into a single mission, and measure both the loss in mission success and the
distributional distance between the delivered handoff state and the successor's
training distribution. The central claim to be tested is that the second quantity
predicts the first. The reinforcement learning implementation, including Proximal
Policy Optimization, the task definitions, the policy sequencer and the evaluation
harness, will be written for this project; the simulator predates it and is treated
as fixed apparatus.
\end{abstract}

\section{Introduction}

Reinforcement learning works by assigning credit to actions from the reward they
eventually produce. That mechanism degrades as the horizon grows. A spacecraft
mission that begins in orbit around one body and ends on the surface of another
spans tens of thousands of simulation steps, and a single terminal reward gives the
agent almost no information about which of those decisions mattered. The action
spaces are not even consistent across the mission: moving a character and flying a
six-degree-of-freedom spacecraft are different control problems joined end to end.

The usual response is temporal abstraction. Rather than learning the whole mission,
the designer decomposes it into skills, trains each against a dense and local reward,
and executes them in sequence. Each skill then faces a tractable problem, and the
resulting system can be inspected one component at a time.

This decomposition carries an assumption that is rarely stated and even more rarely
measured. Each skill is trained on a distribution of starting states chosen by the
designer. When the skills are chained, a skill no longer receives a sampled starting
state; it receives whatever state the preceding skill happened to leave the system in.
If those two distributions disagree, a skill that is individually reliable can fail in
composition without anything being wrong with the skill itself. The failure belongs to
the seam, not to either component.

The analogy is a relay race. Two fast runners do not guarantee a fast relay, because
the baton exchange is its own skill and neither runner practised it.

This project treats that seam as the object of study rather than as an engineering
inconvenience. I will train two control skills independently, chain them, and ask two
questions. First, how much mission success is lost relative to what the individual
success rates predict. Second, whether that loss is explained by a direct measurement
of how far the delivered handoff state falls outside the successor's training
distribution. The second question is the contribution; the first merely establishes
that there is something to explain.

The formal claim to be tested is stated so that it remains informative in either
direction: \emph{mission outcome at a skill handoff is predicted by the overlap
between the delivered state and the successor's training distribution.} If the chain
degrades, the measurement should localise and explain the degradation. If the chain
holds, the measurement should show that the handoff states remained inside the
distribution the successor already knew.

\section{Prior Work}

Temporal abstraction in reinforcement learning is usually formalised through the
options framework of Sutton, Precup and Singh \cite{sutton1999}, which extends the
Markov decision process to temporally extended actions with their own initiation
sets and termination conditions. An option's initiation set is precisely the
statement that a skill is only valid from certain states, which makes the question
of this project a question about whether initiation sets are respected in practice
when skills are trained independently rather than jointly. Later work learns the
decomposition rather than specifying it: the option-critic architecture of Bacon,
Harb and Precup \cite{bacon2017} learns options and their terminations end to end,
and Nachum et al.~\cite{nachum2018} train a higher-level policy that issues goals to
a lower-level one. Both approaches avoid the handoff problem by training the levels
together. This project deliberately keeps the skills independent, because
independently trained skills are the common engineering practice and the case where
the seam is exposed.

The underlying mechanism is covariate shift under a changing state distribution,
which is the same phenomenon that motivates DAgger \cite{ross2011}. There, a policy
cloned from expert demonstrations drifts into states the expert never visited and has
been taught nothing about; the remedy is to collect labels on the states the learner
itself reaches. The structure is identical at a skill handoff, with the predecessor
skill in the role of the thing that produces unfamiliar states, and one of the repairs
proposed below is DAgger applied to the seam rather than to an expert. Domain
randomisation \cite{tobin2017} attacks the same problem from the other side by
widening the training distribution until the deployment distribution falls inside it,
and supplies a second repair to compare against.

The learning algorithm itself is Proximal Policy Optimization \cite{schulman2017},
chosen for stability and for being simple enough to implement correctly from the
paper. Reward shaping follows the potential-based form of Ng, Harada and Russell
\cite{ng1999}, which guarantees that shaping cannot change the optimal policy, a
property that matters here because each skill's reward must be dense enough to train
on without altering what the skill is actually optimising.

What I have not found is direct measurement of the handoff state distribution as a
predictor of composed-system failure. The hierarchical literature largely avoids the
problem by training jointly, and the imitation literature studies distribution shift
between a learner and an expert rather than between two learned skills. That gap is
what this project addresses.

\section{Data}

The data source is a spacecraft simulator built in Godot 4 that predates this project
and is treated as fixed apparatus. No feature work on the simulator is in scope, and
the distinction is stated explicitly so that the contribution being graded is
unambiguous.

The simulator provides analytically propagated planetary orbits, procedurally
generated terrain derived from a per-planet seed, and a rigid-body ship with
six-degree-of-freedom thrust, a finite fuel budget and real gravitational dynamics.
Because the ship obeys gravity while the planets are on rails, the ship's trajectory
is a genuine orbital mechanics problem: orbital elements such as periapsis, apoapsis
and eccentricity are computed from position and velocity rather than scripted.

Data is generated rather than collected. Episodes are produced by stepping 32 ships in
parallel under a headless build, which reaches roughly 5{,}000 environment steps per
second on a laptop CPU, so a three-million-step training run completes in well under
an hour and no GPU is required. One decision is held for six physics ticks, which
places a full episode inside the discount horizon. The experiments below consume on
the order of forty training runs and evaluation sweeps, which is comfortably within a
single machine's budget over four weeks.

Five planets are available, differing in radius, surface gravity and terrain seed.
One is held out of training throughout so that any generalisation claim can be tested
on terrain no policy has seen.

\section{Models}

\textbf{Learning algorithm.} An actor--critic implementation of Proximal Policy
Optimization \cite{schulman2017} written for this project: a clipped surrogate
objective, generalised advantage estimation, separate policy and value heads over a
small fully connected trunk, and running normalisation of the observation vector. The
normaliser's statistics are stored with the exported policy, because they form part of
the policy's definition and are also the instrument used to measure the handoff in
Section~5.

\textbf{Skills.} Two policies are trained independently.

\emph{Orbit insertion} takes a ship approaching a planet on an arbitrary trajectory
and places it in a parking orbit. Success is defined by orbital elements rather than
position: periapsis clear of the terrain, apoapsis inside the planet's capture shell,
eccentricity below a threshold, orbital radius within tolerance of a target, and the
whole condition held long enough that a trajectory merely crossing the band does not
count. Fuel is finite, so the skill is a constrained problem rather than a free one.

\emph{Descent and landing} takes a ship in orbit and lands it intact on a designated
continent, scored on touchdown speed, attitude and whether the landing site is on land.

\textbf{Mission runner.} The component that makes the composition possible, and the
main piece of new engineering. The two skills do not currently share an interface: they
use observation vectors of different width, they disagree about whether the ship's
flight-assist system should be active, and they assume different fuel regimes and
episode horizons. Table~\ref{tab:mismatch} lists the differences. The runner holds the
ship's state across the transition, swaps the active observation builder and policy,
reconfigures the ship, and decides when the handover occurs.

\begin{table}[h]
\centering
\caption{Interface mismatch between the two skills.}
\label{tab:mismatch}
\begin{tabular}{lll}
\toprule
\textbf{Property} & \textbf{Orbit insertion} & \textbf{Descent and landing} \\
\midrule
Observation width & 26 channels & 20 channels \\
Flight assist & Off (it cancels orbital velocity) & On \\
Fuel & Finite budget & Unlimited \\
Episode horizon & 18{,}000 ticks & 27{,}000 ticks \\
\bottomrule
\end{tabular}
\end{table}

\textbf{Baselines.} Each skill is also implemented as a hand-written analytic
controller. These serve two purposes: they establish what the task is worth without
learning, and they provide a control that cannot memorise terrain, which is needed to
interpret any generalisation result.

\section{Experiments}

\textbf{E1: skills in isolation.} Each policy evaluated alone, five random seeds,
200 episodes per seed, actions taken at the distribution mean with exploration noise
disabled. This establishes the reference rates the chain is compared against. The
hand-written controller and a no-input control are scored on the same episodes.

\textbf{E2: the chained mission.} The same policies executed in sequence through the
mission runner, on the same seeds and episode counts. The quantity of interest is the
difference between observed mission success and the product of the individual success
rates from E1. This is the core result and it exists whichever way it comes out.

\textbf{E3: measuring the seam.} At each handover, the delivered state is converted to
a per-channel $z$-score against the successor's stored observation mean and standard
deviation. Two summary statistics are recorded: the largest absolute channel deviation,
and the fraction of channels beyond three standard deviations. These are correlated
against per-episode mission outcome. A relationship between them is the evidence for
the central claim; its absence falsifies the claim and sends the explanation elsewhere.

\textbf{E4: three repairs, compared.} Each repair is applied independently so their
contributions can be separated.

\begin{enumerate}
\itemsep0.15em
\item \emph{Wider initial distribution.} Retrain the successor over a broader range of
starting states, in the spirit of domain randomisation \cite{tobin2017}, so that the
handoff state is no longer unusual.
\item \emph{Handoff fine-tuning.} Run the predecessor, collect the states it actually
produces, and continue training the successor on those states. This is DAgger
\cite{ross2011} applied to a skill boundary rather than to an expert.
\item \emph{Learned switching.} Replace the fixed handover criterion with a small
learned policy that selects the moment of transfer.
\end{enumerate}

\textbf{Controls.} Three controls guard against conclusions that are artefacts rather
than effects, and each is included because it has already caught a false conclusion
during preliminary work on this simulator.

A \emph{frozen-policy control} re-runs any surprising training curve with the learning
rate set to zero. A curve that reproduces with a frozen network is not learning and not
a property of the algorithm. During preliminary work this control converted an apparent
eighty-point training collapse into a diagnosis of drifting observation statistics.

A \emph{baseline-drop control} accompanies any claim that a policy generalises poorly.
The hand-written controller, which has nothing to memorise, is scored on the same
held-out world. If it loses the same amount, the world is harder and no memorisation
has been shown.

\emph{Unit tests on the learning code} cover advantage estimation at truncated episode
boundaries against hand-computed values, a reward-scale assertion that the worst
terminal penalty exceeds the accumulated per-step time cost so that ending an episode
early can never be the cheapest available action, and a check that any imitation-
initialised policy is evaluated on its own before it is trusted.

All results are reported as mean and spread across five seeds on paired episodes. Task
definitions are frozen once E1 is measured; any later change to a reward or tolerance
is reported as an ablation rather than silently folded into the headline result.

\section{Impact}

The practical contribution is a measurement procedure. If handoff distribution overlap
predicts composed-system failure, then anyone chaining independently trained skills has
a cheap diagnostic available before deployment: compare the delivered state against
statistics the policy already stores, and know in advance which seam will fail and
which is safe. That turns a property normally discovered through end-to-end testing
into something checkable at each boundary.

The negative result is also useful. If composition succeeds where the individual
success rates predict it should, that is evidence that independently trained skills
compose more robustly than the covariate-shift argument suggests, and it narrows where
the hierarchical reinforcement learning literature's emphasis on joint training is
actually necessary.

More broadly, autonomous spacecraft operation is a domain where missions are naturally
phased and where each phase already has well-understood classical control laws. Knowing
where learned components can be substituted into such a sequence, and what has to be
verified at the boundaries, is a prerequisite for using them at all.

\section{Timeline}

\begin{table}[h]
\centering
\begin{tabular}{|l|p{10.2cm}|}
\hline
\textbf{Week Of} & \textbf{Description} \\
\hline
October 26 & PPO implementation, environment wrapper and engine bridge; unit tests on the learning code. Gate: landing skill reproduced at its reference rate. \\
\hline
November 2 & Second skill trained; E1 reference measurements at five seeds. Gate: both skills measured and task definitions frozen. \\
\hline
November 9 & Mission runner across the interface mismatch; E2 and E3 run. Gate: chained success rate and seam measurement in hand. \\
\hline
November 16 & E4 repairs implemented and compared; analysis and figures. \\
\hline
November 23 & Finalize paper, code and presentation. \\
\hline
\end{tabular}
\end{table}

The schedule has no slack. Writing the learning algorithm and running the composition
study are each substantial, and four weeks covers both only if nothing is added. If the
second week slips, the fallback is to retain an existing PPO implementation and spend
the remaining time on the study, which preserves the result at the cost of the
implementation exercise. A second fallback reduces the scope to a single handoff rather
than attempting additional skills.

\begin{thebibliography}{9}

\bibitem{bacon2017}
P.-L. Bacon, J. Harb, and D. Precup, ``The option-critic architecture,''
\emph{Proceedings of the AAAI Conference on Artificial Intelligence}, vol.~31,
no.~1, 2017.

\bibitem{nachum2018}
O. Nachum, S. Gu, H. Lee, and S. Levine, ``Data-efficient hierarchical reinforcement
learning,'' \emph{Advances in Neural Information Processing Systems}, vol.~31, 2018.

\bibitem{ng1999}
A. Y. Ng, D. Harada, and S. Russell, ``Policy invariance under reward transformations:
Theory and application to reward shaping,'' \emph{Proceedings of the Sixteenth
International Conference on Machine Learning}, pp.~278--287, 1999.

\bibitem{ross2011}
S. Ross, G. J. Gordon, and J. A. Bagnell, ``A reduction of imitation learning and
structured prediction to no-regret online learning,'' \emph{Proceedings of the
Fourteenth International Conference on Artificial Intelligence and Statistics},
pp.~627--635, 2011.

\bibitem{schulman2017}
J. Schulman, F. Wolski, P. Dhariwal, A. Radford, and O. Klimov, ``Proximal policy
optimization algorithms,'' \emph{arXiv preprint arXiv:1707.06347}, 2017.

\bibitem{sutton1999}
R. S. Sutton, D. Precup, and S. Singh, ``Between MDPs and semi-MDPs: A framework for
temporal abstraction in reinforcement learning,'' \emph{Artificial Intelligence},
vol.~112, no.~1--2, pp.~181--211, 1999.

\bibitem{tobin2017}
J. Tobin, R. Fong, A. Ray, J. Schneider, W. Zaremba, and P. Abbeel, ``Domain
randomization for transferring deep neural networks from simulation to the real
world,'' \emph{IEEE/RSJ International Conference on Intelligent Robots and Systems
(IROS)}, pp.~23--30, 2017.

\end{thebibliography}

\end{document}
